% This is a simple sample document.  For more complicated documents take a look in the exercise tab. Note that everything that comes after a % symbol is treated as comment and ignored when the code is compiled.

\documentclass{article} % \documentclass{} is the first command in any LaTeX code.  It is used to define what kind of document you are creating such as an article or a book, and begins the document preamble

\usepackage{amsmath} % \usepackage is a command that allows you to add functionality to your LaTeX code

\title{Obliczenia Naukowe} % Sets article title
\author{Szymon Pestrowicz} % Sets authors name
\date{\today} % Sets date for date compiled

% The preamble ends with the command \begin{document}
\begin{document} % All begin commands must be paired with an end command somewhere
	\maketitle % creates title using information in preamble (title, author, date)
	
	\section*{Zadanie 1} % creates a section
	    
		\begin{tabular}{|c|c|c|}
		\hline
		 & Lista 1 & Lista 2 \\
		\hline
		Iloczyn skalarny w przód Float32 & -0.4999443 & -0.4999443 \\
		\hline
		Iloczyn skalarny w tył Float32 & -0.4543457 & -0.4543457 \\
		\hline
		Iloczyn skalarny malejaco FLoat32 & -0.5 & -0.5\\
		\hline
		Iloczyn skalarny rosnaco FLoat32 &  -0.5 & -0.5\\
		\hline
		Iloczyn skalarny w przód Float64 & 1.0251881368296672e-10 &-0.004296342739891585\\
		\hline
		Iloczyn skalarny w tył Float64 & -1.5643308870494366e-10 &-0.004296342998713953 \\
		\hline
		Iloczyn skalarny malejaco FLoat64 & 0.0 &-0.004296342842280865 \\
		\hline
		Iloczyn skalarny rosnaco FLoat64 & 0.0 &-0.004296342842280865 \\
		\hline
		\end{tabular}

	\section*{Zadanie 2} % creates a section


	\section*{Zadanie 3} % creates a section

	\section*{Zadanie 4} % creates a section

	\section*{Zadanie 5} % creates a section

	\section*{Zadanie 6} % creates a section	    


\end{document} % This is the end of the document